\section{Introduction}\label{section1}

Non-Hermitian Hamiltonians and complex extension of quantum mechanics have
recently received a lot of attention (see review \cite{[1]}). This f\/ield of
mathematical physics came into f\/lower due to the eigenvalue problem%
\begin{equation}
-\psi ^{\prime \prime }(x)+ix^{3}\psi (x)=E\psi (x),\qquad -\infty
<x<\infty ,  \label{1.1}
\end{equation}%
where $\psi (x)$ is a desired solution, $i=\sqrt{-1}$ is the imaginary unit,
$E$ is a complex parameter (``energy'' or spectral parameter).

A complex value $E_{0}$ of the parameter $E$ is called an \textit{eigenvalue}
of equation (\ref{1.1}) if this equation with $E=E_{0}$ has a nontrivial
(non-identically zero) solution $\psi _{0}(x)$ that belongs to the Hilbert
space $L^{2}(-\infty ,\infty )$ of complex-valued functions $f$ def\/ined on $%
(-\infty ,\infty )$ such that%
\[
\int_{-\infty }^{\infty }\left\vert f(x)\right\vert ^{2}dx<\infty
\]%
with the inner (scalar) product%
\begin{equation}
\left\langle f,g\right\rangle =\int_{-\infty }^{\infty }f(x)\overline{g}(x)dx
\label{1.2}
\end{equation}%
in which the bar over a function denotes the complex conjugate. The function
$\psi _{0}(x)$ is called an \textit{eigenfunction} of equation (\ref{1.1}),
corresponding to the eigenvalue $E_{0}.$

If we def\/ine the operator $S:D\subset L^{2}(-\infty ,\infty )\rightarrow
L^{2}(-\infty ,\infty )$ with the domain $D$ consisting of the functions $%
f\in L^{2}(-\infty ,\infty )$ that are dif\/ferentiable, with the derivative $%
f^{\prime }$ absolutely continuous on each f\/inite subinterval of $(-\infty
,\infty )$ and such that%
\[
-f^{\prime \prime }+ix^{3}f\in L^{2}(-\infty ,\infty ),
\]%
by letting%
\begin{equation}
Sf=-f^{\prime \prime }+ix^{3}f\qquad \text{for} \quad f\in D,  \label{1.3}
\end{equation}%
then eigenvalue problem (\ref{1.1}) can be written as $S\psi =E\psi $, $\psi
\in D,$ $\psi \neq 0.$ Consequently, eigenvalues of equation (\ref{1.1}) are
eigenvalues of the operator~$S$.

Because of the fact that the operator $S$ def\/ined by (\ref{1.3}) has the
complex non-real coef\/f\/icient (potential) equal to $ix^{3},$ this operator is
not Hermitian with respect to the inner product (\ref{1.2}), that is, we
cannot state that%
\[
\left\langle Sf,g\right\rangle =\left\langle f,Sg\right\rangle \qquad \text{for all} \ \ f,g\in D.
\]%
Therefore it is not obvious that the eigenvalues of $S$ may be real
(remember that the eigenvalues of any Hermitian operator are real and the
eigenvectors corresponding to the distinct eigenvalues are orthogonal).

Around 1992 Bessis and Zinn-Justin had noticed on the basis of numerical
work that some of the eigenvalues of equation (\ref{1.1}) seemed to be real
and positive and they conjectured (not in print) that for equation (\ref{1.1}) the eigenvalues are all real and positive.

In 1998, Bender and Boettcher \cite{[2]} generalized the BZJ conjecture,
namely, they conjectured (again on the basis of numerical analysis) that the
eigenvalues of the equation%
\begin{equation}
-\psi ^{\prime \prime }(x)-(ix)^{m}\psi (x)=E\psi (x),\qquad -\infty
<x<\infty ,  \label{1.4}
\end{equation}%
are all real and positive provided $m\geq 2.$ Note that in equation (\ref{1.4}) $m$ is an arbitrary positive real number and equation (\ref{1.1})
corresponds to the choice $m=3$ in (\ref{1.4}).

Bender and Boettcher assumed that the reason for reality of eigenvalues of (%
\ref{1.4}) (in particular of (\ref{1.1})) must be certain symmetry property
of this equation, namely, the so-called $PT$-symmetry of it (for more
details see \cite{[1]}).

$P$ (parity) and $T$ (time reversal) operations are def\/ined by%
\[
Pf(x)=f(-x)\qquad \text{and} \qquad Tf(x)=\overline{f}(x),
\]%
respectively. A Hamiltonian%
\begin{equation}
H=-\frac{d^{2}}{dx^{2}}+V(x)  \label{1.5}
\end{equation}%
with complex potential $V(x)$ is called $PT$-symmetric if it commutes with
the composite operation~$PT$:%
\[
\left[ H,PT\right] =HPT-PTH=0.
\]%
It is easily seen that $PT$-symmetricity of $H$ given by (\ref{1.5}) is
equivalent to the condition%
\[
\overline{V}(-x)=V(x).
\]%
The potentials $ix^{3}$ and $-(ix)^{m}$ of equations (\ref{1.1}) and (\ref%
{1.4}), respectively, are $PT$-symmetric (note that $ix^{3}$ is not $P$%
-symmetric and $T$-symmetric, separately).

The f\/irst rigorous proof of reality and positivity of the eigenvalues of
equation (\ref{1.4}) was given in 2001 by Dorey, Dunning, and Tateo \cite{[3]} (see also \cite{[4]}).

Note that for $0<m<2$ the spectrum of (\ref{1.4}) considered in $%
L^{2}(-\infty ,\infty )$ is also discrete, however in this case only a
f\/inite number of eigenvalues are real and positive, and remaining
eigenvalues (they are of inf\/inite number) are non-real. Besides, if $m=4k$ $%
(k=1,2,\ldots )$ then for any complex value of $E$ all solutions of equation
(\ref{1.4}) belong to $L^{2}(-\infty ,\infty )$ (the so-called Weyl's
limit-circle case holds) and, therefore, all complex values of $E$ are
eigenvalues of equation~(\ref{1.4}) and hence the spectrum is not discrete
for $m=4k$ $(k=1,2,\ldots )$ if we consider the problem in $L^{2}(-\infty
,\infty ).$ In order to have in the case $m\geq 4$ a problem with discrete
spectrum Bender and Boettcher made in~\cite{[2]} an important observation
that the equation can be considered on an appropriately chosen complex
contour. Namely, it is suf\/f\/icient to consider the equation%
\begin{equation}
-\psi ^{\prime \prime }(z)-(iz)^{m}\psi (z)=E\psi (z),\qquad z\in \Gamma
\label{1.6}
\end{equation}%
together with the condition that%
\begin{equation}
\left\vert \psi (z)\right\vert \rightarrow 0 \ \text{ exponentially as }z\text{
moves of\/f to inf\/inity along }\Gamma ,  \label{1.7}
\end{equation}%
where $\Gamma $ is a contour of the form%
\[
\Gamma =\{z=x-i\left\vert x\right\vert \tan \delta :-\infty <x<\infty \}
\]%
with%
\[
\delta =\frac{\pi (m-2)}{2(m+2)}.
\]%
Note that the contour $\Gamma $ forms an angle in the lower complex $z$-plane, of value $\pi -2\delta $ with the vertex at the origin and symmetric
with respect to the imaginary axis.

Next, Mostafazadeh showed in \cite{[5]} that problem (\ref{1.6}), (\ref{1.7}) is equivalent to f\/inding solutions in $L^{2}(-\infty ,\infty )$ of the
problem%
\begin{gather}
-\psi ^{\prime \prime }(x)+\left\vert x\right\vert ^{m}\psi (x)=E\rho
(x)\psi (x),\qquad x\in (-\infty ,0)\cup (0,\infty ),  \label{1.8}
\\
\psi (0^{-})=\psi (0^{+}),\qquad \psi ^{\prime }(0^{-})=e^{2i\delta
}\psi ^{\prime }(0^{+}),  \label{1.9}
\end{gather}%
where%
\begin{equation}
\rho (x)=\left\{
\begin{array}{lll}
e^{2i\delta } & \text{if} & x<0, \\
e^{-2i\delta } & \text{if} & x>0.%
\end{array}%
\right.  \label{1.10}
\end{equation}

The main distinguishing features of problem (\ref{1.8}), (\ref{1.9}) are
that it involves a complex-valued coef\/f\/icient function $\rho (x)$ of the
form (\ref{1.10}) and that transition conditions (impulse conditions) of the
form (\ref{1.9}) are presented which also involve a complex coef\/f\/icient.
Such a problem is non-Hermitian with respect to the usual inner product (\ref%
{1.2}) of space $L^{2}(-\infty ,\infty ).$

Our aim in this paper is to construct and investigate a discrete version of
problem (\ref{1.8}), (\ref{1.9}). Discrete equations (dif\/ference equations)
form a reach f\/ield, both interesting and useful \cite{[6],[7]}. Discrete
equations arise when dif\/ferential equations are solved approximately by
discretization. On the other hand they often arise independently as
mathematical models of many practical events. Discrete equations can easily
been algorithmized to solve them on computers. There is only a small body of
work concerning discrete non-Hermitian quantum systems. Some examples are
\cite{[8],[9],[10],[11],[12],[13],[14],[15]}. Note that in \cite{[15]} the author considered a f\/inite
discrete interval version of the inf\/inite discrete interval problem (\ref{1.12}), (\ref{1.13}) formulated below, and found conditions that ensure the
reality of the eigenvalues. In the f\/inite discrete interval case with the
zero boundary conditions the problem is reduced to the eigenvalue problem
for a f\/inite-dimensional tridiagonal matrix.

Let $
\mathbb{Z}
$ denote the set of all integers. For any $l,m\in
\mathbb{Z}
$ with $l\leq m$, $[l,m]$ will denote the \textit{discrete interval }being
the set $\{l,l+1,\ldots ,m\}.$ Semi-inf\/inite intervals of the form $(-\infty
,l]$ and $[l,\infty )$ will denote the discrete sets $\{\ldots ,l-2,l-1,l\}$
and $\{l,l+1,l+2,\ldots \},$ respectively. Throughout the paper all
intervals will be discrete intervals. Let us set%
\begin{equation}
\mathbb{Z}_{0}=\mathbb{Z}
\setminus \left\{ 0,1\right\} =\left\{ \ldots ,-3,-2,-1\right\} \cup
\left\{ 2,3,4,\ldots \right\} =(-\infty ,-1]\cup \lbrack 2,\infty ).
\label{1.11}
\end{equation}%
We of\/fer a discrete version of problem (\ref{1.8}), (\ref{1.9}) to be
\begin{gather}
-\Delta ^{2}y_{n-1}+q_{n}y_{n}=\lambda \rho _{n}y_{n},\qquad n\in
\mathbb{Z}_{0},  \label{1.12}
\\
y_{-1}=y_{1},\qquad \Delta y_{-1}=e^{2i\delta }\Delta y_{1},
\label{1.13}
\end{gather}%
where $y=\left( y_{n}\right) _{n\in \mathbb{Z}
}$ is a desired solution, $\Delta $ is the forward dif\/ference operator
def\/ined by $\Delta y_{n}=y_{n+1}-y_{n}$ so that $\Delta
^{2}y_{n-1}=y_{n+1}-2y_{n}+y_{n-1},$ the coef\/f\/icients $q_{n}$ are real
numbers given for $n\in \mathbb{Z}_{0},$ $\delta \in \lbrack 0,\pi /2),$ and $\rho _{n}$ are given for $n\in
\mathbb{Z}_{0}$ by
\begin{equation}
\rho _{n}=\left\{
\begin{array}{lll}
e^{2i\delta } & \text{if} & n\leq -1, \\
e^{-2i\delta } & \text{if} & n\geq 2.%
\end{array}%
\right.  \label{1.14}
\end{equation}

One of the main results of the present paper is that if%
\begin{equation}
q_{n}\geq c>0 \qquad \text{for} \ \ n\in \mathbb{Z}_{0}  \label{1.15}
\end{equation}%
and%
\begin{equation}
\lim_{\left\vert n\right\vert \rightarrow \infty }q_{n}=\infty ,
\label{1.16}
\end{equation}%
then the spectrum of problem (\ref{1.12}), (\ref{1.13}) is discrete.

The paper is organized as follows. In Section~\ref{section2}, we choose a
suitable Hilbert space and def\/ine the main linear operators $L$, $M$, and $%
A=M^{-1}L$ related to problem (\ref{1.12}), (\ref{1.13}). Using these
operators we introduce the concept of the spectrum for problem (\ref{1.12}),
(\ref{1.13}). In Section~\ref{section3}, we demonstrate non-Hermiticity of the
operators $L$ and $A.$ In Section~\ref{section4}, we present general properties
of solutions of equations of type (\ref{1.12}), (\ref{1.13}). In Section~\ref{section5}, we construct two special solutions of problem (\ref{1.12}), (\ref{1.13}) under the condition (\ref{1.15}). Using these solutions we show in
Section~\ref{section6} that the operator $L$ is invertible and we describe the
structure of the inverse opera\-tor~$L^{-1}$. Finally, in Section~\ref{section7},
we show that the operator $L^{-1}$ is completely continuous if, in addition,
the condition (\ref{1.16}) is satisf\/ied. This fact yields the discreteness
of the spectrum of problem (\ref{1.12}), (\ref{1.13}).



\section{The concept of the spectrum for problem (\ref{1.12}), (\ref{1.13})}\label{section2}

In order to introduce the concept of the spectrum for problem (\ref{1.12}), (%
\ref{1.13}), def\/ine the Hilbert space $l_{0}^{2}$ of complex sequences $%
y=\left( y_{n}\right) _{n\in
\mathbb{Z}
_{0}}$ such that%
\[
\sum_{n\in
\mathbb{Z}
_{0}}\left\vert y_{n}\right\vert ^{2}<\infty
\]%
with the inner product and norm%
\[
\left\langle y,z\right\rangle =\sum_{n\in
\mathbb{Z}
_{0}}y_{n}\overline{z}_{n},\qquad \left\Vert y\right\Vert =\sqrt{%
\left\langle y,y\right\rangle }=\Big\{ \sum_{n\in
\mathbb{Z}
_{0}}\left\vert y_{n}\right\vert ^{2}\Big\} ^{1/2},
\]%
where $\mathbb{Z}
_{0}$ is def\/ined by (\ref{1.11}) and the bar over a complex number denotes
the complex conjugate.

Now we try to rewrite problem (\ref{1.12}), (\ref{1.13}) in the form of an
equivalent vector equation in~$l_{0}^{2}$ using appropriate operators.
Denote by $D$ the linear set of all vectors $y=\left( y_{n}\right) _{n\in
\mathbb{Z}
_{0}}\in l_{0}^{2}$ such that $\left( q_{n}y_{n}\right) _{n\in \mathbb{Z}
_{0}}\in l_{0}^{2}.$ Taking this set as the domain of $L,$ $L:D\subset
l_{0}^{2}\rightarrow l_{0}^{2}$ is def\/ined by
\[
\left( Ly\right) _{n}=-\Delta ^{2}y_{n-1}+q_{n}y_{n}\qquad \text{for} \ \ n\in
\mathbb{Z}
_{0},
\]%
where $y_{0}$ and $y_{1}$ are def\/ined from the equations
\[
y_{-1}=y_{1},\qquad \Delta y_{-1}=e^{2i\delta }\Delta y_{1}.
\]%
Note that $y_{0}$ and $y_{1}$ are needed when we evaluate $\left( Ly\right)
_{n}$ for $n=-1$ and $n=2,$ respectively:%
\begin{gather*}
(Ly)_{-1}  = -\Delta ^{2}y_{-2}+q_{-1}y_{-1}  =(y_{0}-2y_{-1}+y_{-2})+q_{-1}y_{-1},
\\
(Ly)_{2}  = -\Delta ^{2}y_{1}+q_{2}y_{2}  = (y_{3}-2y_{2}+y_{1})+q_{2}y_{2}
\end{gather*}
and (\ref{1.13}) gives for $y_{1}$, $y_{0}$ the expressions%
\begin{gather*}
y_{1}  = y_{-1}, \\
y_{0}  = y_{-1}+e^{2i\delta }(y_{2}-y_{1})  =\big(1-e^{2i\delta }\big)y_{-1}+e^{2i\delta }y_{2}.
\end{gather*}

Next, def\/ine the operator $M:l_{0}^{2}\rightarrow l_{0}^{2}$ by%
\[
\left( My\right) _{n}=\rho _{n}y_{n}\qquad \text{for} \ \ n\in
\mathbb{Z}_{0},
\]
where $\rho _{n}$ is given by (\ref{1.14}). Obviously, the adjoint $M^{\ast
} $ of $M$ is def\/ined by%
\[
(M^{\ast }y)_{n}=\overline{\rho }_{n}y_{n},\qquad n\in
\mathbb{Z}_{0},
\]%
and since $\left\vert \rho _{n}\right\vert =1,$ we get that $M$ is a unitary
operator:%
\[
MM^{\ast }=M^{\ast }M=I,
\]
where the asterisk denotes the adjoint operator and $I$ is the identity
operator.

Therefore problem (\ref{1.12}), (\ref{1.13}) can be written as
\[
Ly=\lambda My,\qquad y\in D,
\qquad \mbox{or} \qquad
M^{-1}Ly=\lambda y,\qquad y\in D.
\]%
This motivates to introduce the following def\/inition.

\begin{definition}
\label{Def2.1}By the spectrum of problem (\ref{1.12}), (\ref{1.13}) is meant
the spectrum of the operator $A=M^{-1}L$ with the domain $D$ in the space $%
l_{0}^{2}.$
\end{definition}

Remember that (see \cite{[16]}) if $A$ is a linear operator with a domain
dense in a Hilbert space, then a complex number $\lambda $ is called a
\textit{regular point }of the operator $A$ if the inverse $\left( A-\lambda
I\right) ^{-1}$ exists and represents a bounded operator def\/ined on the
whole space. All other points of the complex plane comprise the \textit{spectrum} of the operator $A.$ Obviously the eigenvalues $\lambda $ of an
operator belong to its spectrum, since the operator $\left( A-\lambda
I\right) ^{-1}$ does not exist for such points (the operator $A-\lambda I$
is not one-to-one). The set of all eigenvalues is called the \textit{point
spectrum} of the operator. The spectrum of the operator $A$ is said to be
\textit{discrete }if it consists of a denumerable (i.e., at most countable)
set of eigenvalues with no f\/inite point of accumulation.

A linear operator acting in a Hilbert space and def\/ined on the whole space
is called \textit{completely continuous }if it maps bounded sets into
relatively compact sets (a set is called \textit{relatively compact} if
every inf\/inite subset of this set has a limit point in the space, that may
not belong to the set). Any completely continuous operator is bounded and
hence its spectrum is a compact subset of the complex plane. As is well
known~\cite{[16]}, every nonzero point of the spectrum of a completely
continuous operator is an eigenvalue of f\/inite multiplicity (that is, to
each eigenvalue there correspond only a f\/inite number of linearly
independent eigenvectors); the set of eigenvalues is at most countable and
can have only one accumulation point $\lambda =0.$ It follows that if a
linear operator $A$ with a domain dense in a Hilbert space is invertible and
its inverse $A^{-1}$ is completely continuous, then the spectrum of $A$ is
discrete.

In this paper we show that the operator $L$
 is invertible under the
condition (\ref{1.15}) and that its inverse $L^{-1}$ is a completely
continuous operator if, in addition, the condition (\ref{1.16}) is
satisf\/ied. This implies that the operator $A=M^{-1}L$ is invertible and $%
A^{-1}=L^{-1}M$ is a completely continuous operator. Hence the spectrum of
the operator $A$ is discrete.


\section[Non-Hermiticity of the operators  $L$ and  $A$]{Non-Hermiticity of the operators  $\boldsymbol{L}$ and  $\boldsymbol{A}$}\label{section3}

Let $(f_{k})$ be a given complex sequence, where $k\in
\mathbb{Z}.$ The forward and backward dif\/ference operators $\Delta $ and $\nabla $ are
def\/ined by%
\[
\Delta f_{k}=f_{k+1}-f_{k}\qquad \text{and} \qquad \nabla f_{k}=f_{k}-f_{k-1},
\]%
respectively. We easily see that%
\begin{gather*}
\nabla f_{k}=\Delta f_{k-1},
\\
\Delta ^{2}f_{k}=\Delta (\Delta f_{k})=f_{k+2}-2f_{k+1}+f_{k},
\\
\nabla ^{2}f_{k}=\nabla (\nabla f_{k})=f_{k}-2f_{k-1}+f_{k-2},
\\
\Delta \nabla f_{k}=f_{k+1}-2f_{k}+f_{k-1}=\nabla \Delta f_{k}=\Delta
^{2}f_{k-1}=\nabla ^{2}f_{k+1}.
\end{gather*}
For any integers $a,b\in
\mathbb{Z}
$ with $a<b$ we have the summation by parts formulas%
\begin{gather}
\sum_{k=a}^{b}(\Delta f_{k})g_{k}=f_{k+1}g_{k}\big|_{a-1}^{b}-\sum_{k=a}^{b}f_{k}(\nabla
g_{k})=f_{b+1}g_{b}-f_{a}g_{a-1}-\sum_{k=a}^{b}f_{k}(\nabla g_{k}),
\nonumber\\
\sum_{k=a}^{b}(\nabla f_{k})g_{k}=f_{k}g_{k+1}\big|_{a-1}^{b}-\sum_{k=a}^{b}f_{k}(\Delta
g_{k})=f_{b}g_{b+1}-f_{a-1}g_{a}-\sum_{k=a}^{b}f_{k}(\Delta g_{k}),
\label{3.1}
\\
\sum_{k=a}^{b}(\Delta \nabla f_{k})g_{k}=(\Delta f_{k})g_{k}\big|_{a-1}^{b}-\sum_{k=a}^{b}(\nabla f_{k})(\nabla g_{k}),  \label{3.2}
\\
\sum_{k=a}^{b}(\Delta \nabla f_{k})g_{k}=(\Delta f_{k})g_{k+1}\big|_{a-1}^{b}-\sum_{k=a}^{b}(\Delta f_{k})(\Delta g_{k}),  \nonumber %\label{3.3}
\\
\sum_{k=a}^{b}[(\Delta \nabla f_{k})g_{k}-f_{k}(\Delta \nabla
g_{k})]=[(\Delta f_{k})g_{k}-f_{k}(\Delta g_{k})]_{a-1}^{b}
\nonumber\\
\qquad{}
=[(\Delta f_{b})g_{b}-f_{b}(\Delta g_{b})]-[(\Delta
f_{a-1})g_{a-1}-f_{a-1}(\Delta g_{a-1})].  \label{3.4}
\end{gather}
